\documentclass[14pt, a4paper]{extarticle}


\usepackage{timcw}

\begin{document}

% Титульный лист:
% #1 – тема, #2 – группа, #3 – ФИО, #4 – город и год
\cwtitle{Тема курсового проекта}{Д-Э-??-??}{Ворошилов М.\,В.}{Москва 2026}{09.03.02 Информационные системы и технологии}{Большие данные и машинное обучение}
\newpage

\tableofcontents
\newpage

\intro

Текст введения.

\newpage
\section{Глава 1. Теоретическая часть}
Text \cite{book1}.

Текст текст \( \int_a^b f(x)\,dx. \)
\[
\int_a^b f(x)\,dx.
\]

\subsection{Пункт}
\begin{equation}\label{integral}
  \int_{a}^b \sin x\,dx
\end{equation}
\begin{equation}\label{summa}
  \sum\limits_{i=1}^n x_i
\end{equation}
\begin{equation}\label{alphabeta}
  \alpha + \beta
\end{equation}

\subsection{Пункт}

\newpage
\section{Глава 2. Практическая часть}
text (\ref{summa})

\subsection{Пункт}
text (\ref{alphabeta})

\subsection{Пункт}

\newpage
\section{Название}
\subsection{Пункт}
\subsection{Пункт}

\newpage
\section{Название}
\subsection{Пункт}
\subsection{Пункт}

\newpage
\conclusion

Текст заключения.

\newpage
\centeredrefname
\clearpage
\addcontentsline{toc}{section}{Список использованных источников}
\begin{thebibliography}{99}
\bibitem{book3} Author Book3
\bibitem{book1} Author Book1
\bibitem{book2} Author Book2
\end{thebibliography}

\newpage
\applications

Текст приложений.

\end{document}